% !PRINTED 605
\weberpage{605}
%<<CONT
liebigen Grade $n$, mit reellen Co\"efficienten, die zwei conjugirt imagin\"are und $n - 2$ reelle Wurzeln haben; man kann ja in
\[
f(x) = (x - \alpha_1)(x - \alpha_2) \ldots (x - \alpha_n)
\]
\[
= x^n + a_1 x^{n-1} + a_2 x^{n-2} + \cdots + a_n
\]
$\alpha_1, \alpha_2$ beliebig conjugirt imagin\"ar und die \"ubrigen $\alpha$ reell annehmen. Die Anzahl der reellen und imagin\"aren Wurzeln von $f(x)$ \"andert sich aber nicht, wenn die Co\"efficienten innerhalb gewisser endlicher Grenzen stetig ver\"andert werden (vergl. den siebenten Abschnitt), und folglich giebt es auch solche Gleichungen, deren Co\"efficienten rationale Zahlen sind, da man in beliebiger N\"ahe von irgend welchen gegebenen reellen Zahlen immer rationale Zahlen finden kann.

Es ist also nur noch nachzuweisen, dass man diese rationalen Co\"efficienten so w\"ahlen kann, dass $f(x)$ irreducibel wird. Dies ergiebt sich aber aus folgendem Satze:

3. \gesperrt{Ist $p$ eine Primzahl, $c_0, c_1, c_2 \ldots c_n$ eine Reihe ganzer Zahlen, von denen $c_0$ durch $p$ nicht theilbar, $c_1, c_2 \ldots c_n$ durch $p$ theilbar, aber $c_n$ nicht durch $p^2$ theilbar ist, so ist die Function}
\[
\varphi(x) = c_0 x^n + c_1 x^{n-1} + \cdots + c_n
\]
\gesperrt{irreducibel.}

Denn wenn $\varphi(x)$ in zwei Factoren mit rationalen Co\"efficienten zerf\"allt, so k\"onnen (nach \S.\,2) die Factoren auch ganzzahlig angenommen werden. Sei also
\[
\varphi(x) = (\alpha_0 x^h + \alpha_1 x^{h-1} + \cdots + \alpha_h)(\beta_0 x^k + \beta_1 x^{k-1} + \cdots + \beta_k),
\]
$h$ und $k$ beide gr\"osser als Null und ihre Summe gleich $n$.

Da $\alpha_h \beta_k = c_n$ ist, so muss einer der beiden Factoren $\alpha_h, \beta_k$ durch $p$ theilbar sein, der andere nicht. Es m\"oge $\beta_k$ durch $p$ theilbar, $\alpha_h$ nicht theilbar sein, und da nicht alle $\beta$ durch $p$ theilbar sein k\"onnen, weil sonst auch $c_0$ durch $p$ theilbar w\"are, so sei $\beta_\nu$ nicht durch $p$ theilbar, $\beta_{\nu+1}, \beta_{\nu+2} \ldots \beta_k$ durch $p$ theilbar. Der Co\"efficient von $x^{k-\nu}$ in dem Product der beiden Factoren ist dann $\alpha_h \beta_\nu + \alpha_{h-1}\beta_{\nu+1} + \cdots$, also durch $p$ nicht theilbar. Es m\"usste also $k - \nu = n$ sein, was nicht m\"oglich ist, da schon $k < n$ sein muss.

Setzt man also f\"ur die Co\"efficienten von $f(x)$:
\[
a_1 = \frac{c_1}{c_0}, \quad a_2 = \frac{c_2}{c_0} \cdots a_n = \frac{c_n}{c_0},
\]
% [page alignment corrected 2026-08-07] Printed 605 ends immediately after
% the coefficient display above.  The following theorem/prose belongs to
% printed 606 and has been moved to i0631.
%CONT>>
